% !PRINTED 613
\weberpage{613}
%<<CONT

4. \gesperrt{The number of non-equivalent irrational numbers of the discriminant $\varDelta$ is equal to the number of classes of forms of the field $\Omega$}\footnote{According to the remark made at the close of \S.\,128 of the first volume on the \gesperrt{Gauss} theory of quadratic forms, the class number of the quadratic forms according to the \gesperrt{Gauss} definition does not always agree with the number of ideal classes of the field $\Omega$. To establish the agreement, one must in some cases divide the ideal classes further into two classes each. In our theory of the equivalence of the numbers $\omega$, where proper equivalence is not distinguished from improper, this further division is not required. Cf. \gesperrt{Dirichlet-Dedekind}, Zahlentheorie. 4. Auflage, S. 578, 584.}\gesperrt{.}

\begin{center}
§. 168.

\gesperrt{Composition of the quadratic irrational numbers.}
\end{center}

We have seen in \S.\,166 that for every integral functional $\varphi$ in the field $\Omega$ a certain quadratic irrational number $\omega$ of the discriminant $\varDelta$ can be found.

But we must now still investigate how far this number $\omega$ is determined by the field $\Omega$ and by the functional $\varphi$.

By \S.\,166, (1), the basis of $\varphi$

\weq{1}{\alpha_1 = a_{1,1},\quad \alpha_2 = a_{1,2}+a_{2,2}\Theta}

is defined by the condition that $a_{1,1}$ is the least positive integral rational number that is divisible by $\varphi$, and $a_{2,2}$ is then determined by

\weq{2}{N_a(\varphi) = a_{1,1}a_{2,2}}

and thus the two numbers $a_{1,1},a_{2,2}$ are uniquely determined by $\varphi$. It remains to ask how far $a_{1,2}$ is also determined. But the number $a_{1,2}$ is given, with respect to the modulus $\varphi$, by the congruence $a_{1,2}+a_{2,2}\Theta\equiv 0$, and since $a_{1,2}$ must also be an integral rational number, $a_{1,2}$ is determined up to a multiple of $a_{1,1}$. This multiple of $a_{1,1}$ remains in fact arbitrary, since, if $\alpha_1,\alpha_2$ is a basis of $\varphi$, the same holds, for every integral rational $k$, of $\alpha_1,\alpha_2+k\alpha_1$.

Now since $\alpha_2:\alpha_1=\omega$ is the quadratic irrational number sought, we obtain the theorem:
%CONT>>
